\documentclass[conference]{IEEEtran}
\IEEEoverridecommandlockouts

% ── Paketler ────────────────────────────────────────────────────────────────
\usepackage{cite}
\usepackage{amsmath,amssymb,amsfonts}
\usepackage{graphicx}
\usepackage{textcomp}
\usepackage{xcolor}
\usepackage{booktabs}
\usepackage{multirow}
\usepackage{url}
\usepackage{hyperref}

% ── Başlık ve Yazarlar ───────────────────────────────────────────────────────
\title{Privacy-Preserving Traffic Anomaly Detection with Federated Learning:\\
A Mininet/SDN Emulation and Ablation Study on PeMS04}

\author{
  \IEEEauthorblockN{Enes Avcı}
  \IEEEauthorblockA{
    \textit{Department of Computer Engineering}\\
    \textit{Bursa Uludağ University}\\
    Bursa, Turkey\\
    % enes.avci@uludag.edu.tr
  }
  \and
  \IEEEauthorblockN{Murtaza Cicioğlu}
  \IEEEauthorblockA{
    \textit{Department of Computer Engineering}\\
    \textit{Bursa Uludağ University}\\
    Bursa, Turkey\\
    % murtaza.cicioglu@uludag.edu.tr
  }
}

\begin{document}
\maketitle

% ── Abstract ────────────────────────────────────────────────────────────────
\begin{abstract}
Intelligent Transportation Systems (ITS) generate continuous streams of
sensor data that are conventionally aggregated at a centralised Traffic
Management Centre (TMC), posing significant privacy risks and bandwidth
demands.
In this paper, we propose a privacy-preserving traffic anomaly detection
framework based on Federated Learning (FL) and evaluate it on the
real-world PeMS04 dataset (307 sensors, 59 days, 5-minute intervals).
We emulate a four-node ITS topology using Mininet and the os-ken
OpenFlow~1.3 SDN controller, and implement the FedAvg aggregation
algorithm with an Isolation Forest (IF) anomaly detector.
A systematic three-scenario ablation study is conducted: (S1)~centralised
baseline, (S2)~IID-distributed FL, and (S3)~Non-IID FL with
heterogeneous sensor profiles.
Experimental results show that S2 achieves an average F1 score of
$0.703$, only $8.9\%$ below the centralised baseline ($0.772$), while
keeping raw sensor data local.
S3 reveals a $24.3\%$ drop ($0.584$) caused by client drift under
heterogeneous data distributions.
Server-to-client communication is reduced to $0.15$~KB per round through
Federated Threshold Calibration, validating the communication-efficiency
argument of FL-based anomaly detection.
\end{abstract}

\begin{IEEEkeywords}
Federated Learning, Anomaly Detection, Intelligent Transportation
Systems, Software-Defined Networking, Mininet, PeMS04, FedAvg,
Isolation Forest, Non-IID, Client Drift
\end{IEEEkeywords}

% ════════════════════════════════════════════════════════════════════════════
\section{Introduction}
% ════════════════════════════════════════════════════════════════════════════

Modern ITS infrastructures rely on dense networks of IoT sensors—loop
detectors, speed radars, and cameras—that continuously monitor traffic
flow, speed, and occupancy.
Timely detection of anomalies such as accidents, congestion, or sensor
faults is critical for traffic management.
The dominant paradigm transmits raw sensor readings to a centralised TMC,
which introduces three fundamental vulnerabilities:
(i)~\textit{privacy risk}—vehicle trajectory data constitute personally
identifiable information regulated by frameworks such as GDPR and the
Turkish KVKK;
(ii)~\textit{bandwidth pressure}—hundreds of high-frequency sensors
saturate shared communication infrastructure;
(iii)~\textit{data island problem}—sensors owned by different
administrative entities cannot share raw data for joint model training
\cite{liu2021deep}.

Federated Learning (FL) \cite{mcmahan2017communication} addresses these
challenges by training models locally on each client and transmitting
only model parameters to a central aggregator, so raw data never leaves
the sensor node.
While FL has been applied to ITS traffic prediction
\cite{xia2023short,yuan2023fedSTN}, systematic ablation studies that
isolate the contribution of FL under realistic network conditions—using
actual SDN-controlled topologies and TCP sockets—remain scarce.

This paper makes the following contributions:
\begin{itemize}
  \item We implement and evaluate FedAvg-based anomaly detection on a
        Mininet/os-ken SDN emulation with four heterogeneous PeMS04
        sensor nodes, providing a reproducible network-layer testbed.
  \item We conduct a three-scenario ablation study (S1~centralised,
        S2~IID FL, S3~Non-IID FL) that isolates the effects of
        federation and data heterogeneity on F1 score, communication
        cost, and convergence.
  \item We propose \textit{Federated Threshold Calibration}, a
        lightweight aggregation scheme for Isolation Forest that
        reduces server-to-client traffic to $0.15$~KB per round.
\end{itemize}

% ════════════════════════════════════════════════════════════════════════════
\section{Related Work}
% ════════════════════════════════════════════════════════════════════════════

\subsection{FL for ITS Applications}

McMahan et al.~\cite{mcmahan2017communication} introduced FedAvg, which
minimises communication cost by performing multiple local SGD steps
before aggregation.
Li et al.~\cite{li2020federated} systematically characterise the
challenges of statistical and system heterogeneity in FL, motivating our
Non-IID scenario.

In the ITS domain, Xia et al.~\cite{xia2023short} and
Yuan et al.~\cite{yuan2023fedSTN} apply graph-based FL to traffic flow
prediction on PeMS datasets, establishing benchmark performance figures
that we use for dataset validation.
Chellapandi et al.~\cite{chellapandi2024federated} survey FL for
connected vehicles, identifying communication cost and Non-IID data as
open challenges.

\subsection{FL for Anomaly Detection}

Liu et al.~\cite{liu2021deep} propose a communication-efficient
on-device FL framework for time-series anomaly detection in Industrial
IoT, which serves as the direct methodological template of this work.
Tham et al.~\cite{tham2023federated} compare FedAvg, FedProx, and
FedAvg-Adam for anomaly detection in vehicular networks under IID and
Non-IID settings, confirming that Non-IID distributions consistently
degrade global model performance.

\subsection{Research Gap}

Existing FL+ITS works predominantly simulate FL on a single machine
using frameworks such as Flower or PySyft, bypassing the network layer.
Our work explicitly emulates the communication topology with Mininet and
os-ken, measuring real TCP-socket byte counts and enabling future
integration of network impairment (delay, packet loss) without code
changes.

% ════════════════════════════════════════════════════════════════════════════
\section{System Design and Dataset}
% ════════════════════════════════════════════════════════════════════════════

\subsection{Network Topology}

We emulate a four-node ITS network using Mininet~2.3.0
(Linux Bridge mode) and os-ken~2.3.1 as the OpenFlow~1.3 SDN
controller—the actively-maintained OpenStack fork of Ryu
\cite{osken2024}.
Each host ($h_1$–$h_4$, IP 10.0.0.1–4) represents an ITS sensor
acting as an FL client; a fifth host (10.0.0.10) represents the TMC
aggregation server.
All hosts connect to a single OVS switch managed by os-ken via
\texttt{simple\_switch\_13}.
Connectivity is verified with \texttt{pingall} (0\% packet loss,
20/20 received); first-hop RTT~$\approx$1~ms (Packet-In + Flow-Mod
latency), subsequent RTT~$\approx$0.05~ms (flow-table hit).
FL messages are exchanged over TCP sockets with length-prefixed framing
(\texttt{struct.pack('>I', len)}) to handle large payloads safely.

\subsection{Dataset: PeMS04}

\begin{figure}[t]
  \centering
  \includegraphics[width=\columnwidth]{fig06_selected_sensors_weekly.png}
  \caption{Weekly traffic flow patterns of the four selected PeMS04
           sensors (P25/P50/P75/P95). Distinct diurnal and load profiles
           confirm Non-IID data heterogeneity across clients.}
  \label{fig:sensor_profiles}
\end{figure}

The PeMS04 dataset \cite{caltrans2024pems} comprises traffic
measurements from 307 sensors across 29 roads in the San Francisco Bay
Area over 59 days (January–February 2018) at 5-minute intervals,
yielding $16{,}992 \times 307 \times 3$ samples (flow, speed,
occupancy).
No missing values were found.

Four sensors are selected via a percentile-based strategy
(P25, P50, P75, P95 of mean flow) to ensure heterogeneous Non-IID
profiles representative of low-density rural, arterial, on-ramp, and
high-density highway segments, respectively
(sensor IDs: 278, 240, 71, 298).
A notable finding is that flow values never exceed $|z|>3$; anomalies
are therefore driven exclusively by speed and occupancy, which is
consistent with the distributional properties of PeMS data and is
reported as a dataset observation.

\subsection{Anomaly Labelling}

Anomaly labels are generated automatically using Z-score thresholding
($|z|>3$) applied per sensor per feature (OR rule: any feature
exceeding the threshold flags the time step as anomalous).
Z-score statistics ($\mu$, $\sigma$) are computed exclusively on the
training split to prevent data leakage.
A chronological 70/15/15 train/validation/test split is applied
(no shuffling) to preserve temporal integrity.
The resulting contamination rate across sensors ranges from
$1.49\%$ to $1.84\%$, confirming label consistency.

% ════════════════════════════════════════════════════════════════════════════
\section{Methodology}
% ════════════════════════════════════════════════════════════════════════════

\subsection{Ablation Scenarios}

Three scenarios isolate the contribution of each system component:

\begin{itemize}
  \item \textbf{S1 — Centralised Baseline:} All four sensors' data are
        pooled at the TMC. An Isolation Forest is trained on the combined
        dataset. Raw data ($\approx$21~MB) is transmitted to the TMC.
  \item \textbf{S2 — IID FL:} Training data is shuffled and partitioned
        equally across four clients before being assigned to sensor nodes.
        FedAvg is executed over three communication rounds.
  \item \textbf{S3 — Non-IID FL:} Each client retains its own sensor's
        natural traffic distribution (heterogeneous flow profiles),
        reflecting a realistic deployment where data islands are preserved.
\end{itemize}

\subsection{Federated Threshold Calibration}

Standard FL aggregation of tree-based models (e.g., averaging decision
trees) is impractical due to structural incompatibility after
serialisation.
We propose \textit{Federated Threshold Calibration}: each client trains
a full local Isolation Forest and transmits only its anomaly decision
threshold (\texttt{offset\_}) to the TMC.
The server computes a sample-weighted average of client thresholds and
broadcasts the global threshold back to all clients, which update their
local decision boundaries without retraining.
This approach is directly aligned with the communication-efficient FL
philosophy of \cite{liu2021deep}, reducing server-to-client traffic to
$0.15$~KB per round (versus $\approx$1.7~MB for full model transmission).

\subsection{FedAvg Protocol}

\begin{equation}
\hat{\tau}_{\text{global}} = \frac{\sum_{k=1}^{K} n_k \cdot \tau_k}{\sum_{k=1}^{K} n_k}
\label{eq:fedavg_threshold}
\end{equation}

\noindent
where $K=4$ is the number of clients, $n_k$ is the number of training
samples at client $k$, and $\tau_k$ is the local anomaly threshold.
Three aggregation rounds are executed; all four clients participate in
every round (full participation).

\subsection{Evaluation Metrics}

Performance is assessed on the held-out test set (15\% of each sensor's
data) using:
(1)~F1 score (primary metric, harmonic mean of Precision and Recall);
(2)~communication cost (total bytes transmitted over TCP sockets,
measured by the Mininet host processes);
(3)~client drift, defined as the signed difference between global and
local F1 per sensor ($\Delta_k = \text{F1}_{\text{global},k} -
\text{F1}_{\text{local},k}$).

% ════════════════════════════════════════════════════════════════════════════
\section{Experimental Results}
% ════════════════════════════════════════════════════════════════════════════

\subsection{S1: Centralised Baseline}

\begin{figure}[t]
  \centering
  \includegraphics[width=0.75\columnwidth]{fig08_s1_confusion_matrix.png}
  \caption{Confusion matrix of the S1 centralised Isolation Forest
           on the test set (TP=184, FP=2, FN=107, TN=9907).
           High Precision ($0.989$) with moderate Recall ($0.632$)
           indicates conservative threshold behaviour.}
  \label{fig:confusion}
\end{figure}

The centralised Isolation Forest achieves
$\text{F1}=0.7715$, $\text{Precision}=0.9892$, $\text{Recall}=0.6323$
on the test set (Table~\ref{tab:ablation}).
The high Precision–low Recall profile indicates conservative threshold
behaviour: the model raises very few false alarms but misses $36.8\%$
of true anomalies.
This pattern is consistent with the default contamination-based
thresholding of Isolation Forest and serves as the reference upper bound
for the FL scenarios.

\subsection{S2 and S3: Federated Learning Results}

\begin{figure}[t]
  \centering
  \includegraphics[width=\columnwidth]{fig10_ablation_f1.png}
  \caption{Ablation study results: average F1 score across three
           scenarios. S2 (IID FL) achieves $91.1\%$ of centralised
           accuracy; S3 (Non-IID FL) shows a $24.3\%$ drop due to
           client drift.}
  \label{fig:ablation}
\end{figure}

\begin{figure*}[t]
  \centering
  \includegraphics[width=0.82\textwidth]{fig11_client_drift.png}
  \caption{Client drift analysis: local (pre-aggregation) vs.\ global
           (federated) F1 score per sensor under IID (left) and Non-IID
           (right) data distributions. Positive $\Delta$ indicates
           federation improves the client; negative $\Delta$ indicates
           the global threshold degrades local performance.}
  \label{fig:drift}
\end{figure*}

Table~\ref{tab:ablation} summarises the three-scenario comparison.
S2 (IID FL) reaches an average global F1 of $0.703$, a modest
$8.9\%$ reduction relative to S1, while raw data never leaves any
client node.
This trade-off is considered acceptable for privacy-sensitive ITS
deployments.

S3 (Non-IID FL) yields $\text{F1}=0.584$, a $24.3\%$ drop from S1,
attributable to client drift: heterogeneous training distributions
produce divergent local thresholds, and the weighted average degrades
performance for sensors whose traffic profiles differ most from the
global mean.
Sensor~240 exhibits the clearest drift instance
($\Delta_{240} = -0.036$), where the aggregated threshold is pulled
away from its local optimum by the three other sensors.

\subsection{Communication Cost}

\begin{figure}[t]
  \centering
  \includegraphics[width=\columnwidth]{fig12_comm_cost.png}
  \caption{Communication cost comparison across scenarios.
           S1 transmits raw sensor data ($\approx$21~MB) to the TMC.
           FL scenarios transmit only model parameters and a
           $0.15$~KB global threshold per round, keeping raw data local.}
  \label{fig:commcost}
\end{figure}

Table~\ref{tab:comm} compares data transmission across scenarios.
Federated Threshold Calibration reduces server-to-client traffic to
$0.15$~KB per round (total $1.8$~KB over three rounds for four clients),
confirming the communication-efficiency claim of
\cite{liu2021deep}.
Client-to-server traffic ($\approx$1.7~MB per client per round) is
dominated by the serialised Isolation Forest model; transmitting only
the threshold scalar would reduce this to $<$1~KB, which we identify as
a direction for future optimisation.

% ── Ablasyon Tablosu ────────────────────────────────────────────────────────
\begin{table}[t]
\caption{Three-Scenario Ablation: Anomaly Detection Performance (PeMS04, 4 Sensors)}
\label{tab:ablation}
\centering
\begin{tabular}{lccc}
\toprule
\textbf{Scenario} & \textbf{Avg.\ F1} & \textbf{$\Delta$ vs S1} & \textbf{Data Privacy} \\
\midrule
S1 Centralised    & 0.7715 & —          & \texttimes\ (raw data at TMC) \\
S2 IID FL         & 0.7028 & $-8.9\%$   & \checkmark\ (local only) \\
S3 Non-IID FL     & 0.5843 & $-24.3\%$  & \checkmark\ (local only) \\
\bottomrule
\end{tabular}
\end{table}

% ── İletişim Maliyeti Tablosu ────────────────────────────────────────────────
\begin{table}[t]
\caption{Communication Cost Comparison (4 Clients $\times$ 3 FL Rounds)}
\label{tab:comm}
\centering
\begin{tabular}{lcc}
\toprule
\textbf{Scenario} & \textbf{Client$\to$Server} & \textbf{Server$\to$Client} \\
\midrule
S1 Centralised & $\approx$21 MB (raw data) & — \\
S2 IID FL      & $\approx$20.4 MB (model)  & 1.8 KB (threshold) \\
S3 Non-IID FL  & $\approx$20.4 MB (model)  & 1.8 KB (threshold) \\
\bottomrule
\end{tabular}
\end{table}

% ════════════════════════════════════════════════════════════════════════════
\section{Discussion}
% ════════════════════════════════════════════════════════════════════════════

\textbf{Privacy--accuracy trade-off (AS1):}
S2 demonstrates that FL achieves $91.1\%$ of the centralised F1
without ever transmitting raw sensor data, directly answering AS1.
The $8.9\%$ gap is attributable to the aggregation of threshold scalars
rather than full gradient updates, and to the three-round convergence
limit.

\textbf{Communication cost (AS2):}
Server-to-client overhead ($1.8$~KB total) is negligible compared to
the $21$~MB centralised raw-data transfer, validating FL's communication
advantage (AS2).
Client-to-server cost remains high due to full model serialisation;
transmitting only \texttt{offset\_} would reduce this by four orders of
magnitude.

\textbf{Non-IID effect and client drift (AS3):}
The $15.4$~percentage-point gap between S2 and S3 confirms that data
heterogeneity is the dominant factor degrading FL performance (AS3).
The weighted-average aggregation of FedAvg is insufficient to reconcile
divergent threshold distributions; adaptive aggregation strategies
(e.g., FedProx, personalised FL) are reserved for the journal extension.

\textbf{Limitations:}
(i)~Isolation Forest threshold aggregation is a simplification of true
model federation; deep learning models would enable gradient-based
FedAvg.
(ii)~The emulation does not introduce artificial network impairments
(delay, jitter, packet loss); this is planned for the Q1 journal version.
(iii)~Only one benchmark dataset (PeMS04) is used; cross-validation on
PeMS08 and METR-LA is deferred.

% ════════════════════════════════════════════════════════════════════════════
\section{Conclusion}
% ════════════════════════════════════════════════════════════════════════════

We presented a privacy-preserving FL framework for traffic anomaly
detection, evaluated on PeMS04 via a Mininet/os-ken SDN emulation.
Our three-scenario ablation study shows that IID FL achieves $91.1\%$
of centralised accuracy while keeping raw data local, and that Non-IID
data heterogeneity causes a $24.3\%$ F1 degradation through client drift.
Federated Threshold Calibration reduces server-to-client traffic to
$0.15$~KB per round.
These findings provide a reproducible, network-layer baseline for future
FL+ITS research.
Future work will extend the framework with FedProx under realistic
4G/LTE channel models, Differential Privacy, and gradient compression.

% ── Kaynakça ────────────────────────────────────────────────────────────────
\begin{thebibliography}{00}

\bibitem{mcmahan2017communication}
H.~B. McMahan, E.~Moore, D.~Ramage, S.~Hampson, and B.~A. y~Arcas,
``Communication-efficient learning of deep networks from decentralized data,''
in \textit{Proc. 20th Int. Conf. Artif. Intell. Statist. (AISTATS)},
Fort Lauderdale, FL, USA, Apr. 2017, pp. 1273--1282.

\bibitem{li2020federated}
T.~Li, A.~K. Sahu, A.~Talwalkar, and V.~Smith,
``Federated learning: Challenges, methods, and future directions,''
\textit{IEEE Signal Process. Mag.}, vol.~37, no.~3, pp. 50--60, May 2020.

\bibitem{liu2021deep}
Y.~Liu, S.~Garg, J.~Nie, Y.~Zhang, Z.~Xiong, J.~Kang, and M.~S. Hossain,
``Deep anomaly detection for time-series data in industrial IoT:
A communication-efficient on-device federated learning approach,''
\textit{IEEE Internet Things J.}, vol.~8, no.~8, pp. 6348--6358, Apr. 2021.

\bibitem{xia2023short}
M.~Xia, D.~Jin, and J.~Chen,
``Short-term traffic flow prediction based on graph convolutional networks
and federated learning,''
\textit{IEEE Trans. Intell. Transp. Syst.}, vol.~24, no.~1,
pp. 1191--1203, Jan. 2023.

\bibitem{yuan2023fedSTN}
X.~Yuan et al.,
``FedSTN: Graph representation driven federated learning for edge computing
enabled urban traffic flow prediction,''
\textit{IEEE Trans. Intell. Transp. Syst.}, vol.~24, no.~8,
pp. 8738--8748, Aug. 2023.

\bibitem{chellapandi2024federated}
V.~P. Chellapandi, L.~Yuan, C.~G. Brinton, S.~H. Żak, and Z.~Wang,
``Federated learning for connected and automated vehicles: A survey of
existing approaches and challenges,''
\textit{IEEE Trans. Intell. Veh.}, vol.~9, no.~1, pp. 119--137, Jan. 2024.

\bibitem{tham2023federated}
K.~P. Tham et al.,
``Federated learning for anomaly detection in vehicular networks,''
in \textit{Proc. IEEE Veh. Technol. Conf. (VTC-Spring)}, 2023.

\bibitem{osken2024}
OpenStack Foundation,
``os-ken: A component-based SDN framework (actively-maintained fork of Ryu),''
\url{https://opendev.org/openstack/os-ken}, 2024.

\bibitem{caltrans2024pems}
California Dept.\ of Transportation (Caltrans),
``PeMS: Performance measurement system,''
\url{https://pems.dot.ca.gov/}, 2024.

\end{thebibliography}

\end{document}
